\documentclass[11pt,a4paper]{article}
\usepackage[margin=25mm]{geometry}
\usepackage{amsmath,amssymb,amsthm,booktabs,hyperref}
\newtheorem{theorem}{Theorem}\newtheorem{lemma}[theorem]{Lemma}
\newcommand{\Aut}{\operatorname{Aut}}
\title{The order-eleven column of the O'Nan map count}
\author{Human Compression project}\date{}
\begin{document}\maketitle
\begin{abstract}
A pair generates a unique subgroup. Using that partition, we recover all
sixteen entries involving order eleven from three intrinsic subgroup
counts and their class lengths. Only one overlap type, $L_2(11)$, occurs.
The calculation uses the previously retained character vectors, one
additional character table and four rows of the published subgroup atlas.
Nine entries were not determined by the earlier components.
\end{abstract}
\section{Partition by the generated subgroup}
For a finite group $H$ and $m>2$, write
\[
 X_H=\{(r,s):|r|=m,\ |s|=11,\ |rs|=2\},\qquad
 E_H=\{(r,s)\in X_H:\langle r,s\rangle=H\},
\]
and put $x_H=|X_H|$, $e_H=|E_H|$. Orders are intrinsic to the
subgroup, so the counts do not depend on its ambient embedding.
The following is the generated-subgroup partition from our common
counting theorem, used without a full ambient incidence matrix.
\begin{lemma}\label{partition}
Suppose $H_1,\ldots,H_k$ represent every conjugacy class of proper
subgroups that can be generated by a pair in $X_G$, and let $c_i$
be the length of the conjugacy class of $H_i$. Then
\[
 e_G=x_G-\sum_{i=1}^k c_i e_{H_i}.
\]
\end{lemma}
\begin{proof}
Every nongenerating pair generates exactly one proper subgroup $K$.
The sets $E_K$ for distinct actual subgroups are disjoint. Conjugation
gives bijections among the $E_K$ in each conjugacy class, so that class
contributes $c_i e_{H_i}$. Subtract their sum from $x_G$.
\end{proof}
In particular, one must subtract intrinsic \emph{generating}-pair
counts at this stage. Subtracting all pairs in each maximal subgroup
would count pairs in overlaps more than once.

\section{The proper possibilities}
Let $G=O'N$, and put $L=L_2(11)$, $J=J_1$ and $M=M_{11}$.
We retain the complete maximal-subgroup list already used in the
preceding bounds. Exactly three maximal classes have order divisible
by eleven: one of type $J$ and two of type $M$.
The following smaller subgroup facts are the exact portions of the
published atlas~\cite{atlas} used here.
\begin{center}\begin{tabular}{rrrr}\toprule
atlas row&type&order&class length in $G$\\\midrule
6&$J$&175560&2624832\\
23,24&$M$&7920&58183776 each\\
119&$L$&660&698205312\\\bottomrule
\end{tabular}\end{center}
Within $J$, the only maximal types with order divisible by eleven
are $L$ and $11:10$. There are 266 conjugates of $L$ in $J$.
Within $M$, the only such maximal type is $L$, with twelve conjugates.
The two $G$-classes of $M$ are kept separately.
All these $L$ subgroups belong to the single $G$-class in row 119.
The atlas label $11:5:2$ in row 286 supplies a normal subgroup
of order eleven. The retained $J_1$ character table has no element
of order 22; these are the only structural properties of this
order-110 group needed below.
These classifications and class lengths are borrowed inputs;
their coverage is not inferred from a few character identities.

\begin{lemma}\label{L}
Every pair in $X_L$ generates $L$. No pair in $X_{11:10}$ exists.
\end{lemma}
\begin{proof}
The previously proved prime-field lemma says that a proper subgroup
of $L_2(p)$ containing an element of order $p$ lies in a point
stabilizer. For $p=11$ that stabilizer has order $11(11-1)/2=55$,
which is odd. A pair in $X_L$ generates a subgroup containing both
its order-eleven second element and its involutory product, so that
subgroup cannot be proper.

For the second assertion put $b=rs$. The subgroup $P$ of order eleven
is normal. An involution centralizing $P$ would give an element of
order 22, which the containing group $J_1$ does not have. Hence
every involution acts on $P$ by inversion. Since $s\in P$ and $b$ is an involution,
$r=bs^{-1}$ is an involution as well, contrary to $m>2$.
Equivalently one may use the cyclic-normal-subgroup lemma already
proved for the order-nineteen column.
\end{proof}
It follows that the only proper generated subgroups in $J$ or $M$
are copies of $L$. Different copies cannot share such a pair, since
that pair generates its copy. Therefore
\[
 e_J=x_J-266x_L,\qquad e_M=x_M-12x_L. \tag{1}
\]
Every proper generated subgroup in $G$ is consequently a maximal
$J$, a maximal $M$, or an $L$. Thus no unlisted subgroup incidence is
needed. In particular we do not multiply cover weights along chains.

\begin{theorem}\label{formula}
For every $m>2$,
\[
 e_G(m,11)=x_G(m,11)-\frac{|G|}{175560}x_J(m,11)
 -2\frac{|G|}{7920}x_M(m,11)
 +2\frac{|G|}{660}x_L(m,11). \tag{2}
\]
\end{theorem}
\begin{proof}
Apply Lemma~\ref{partition} to the three maximal classes and the
single class of $L$. Their lengths are $|G|/175560$,
$|G|/7920$, $|G|/7920$ and $|G|/660$.
Substitute (1), and use $266\cdot660=175560$ and
$12\cdot660=7920$. The coefficient of $x_L$ is
$(1+2-1)|G|/660=2|G|/660$, as in (2).
\end{proof}

\section{Exact arithmetic and attribution}
For each retained table we use the already established character formula
\[
 x_H(m,11)=\frac1{|H|}\sum_{\chi\in\operatorname{Irr}(H)}
 \frac{v_{\chi,2}v_{\chi,m}v_{\chi,11}}{\chi(1)},
 \qquad v_{\chi,t}=\sum_{|c|=t}|c^H|\chi(c).
\]
The only new table export is $L_2(11)$; its eight vectors are integral.
The reader performs exact rational arithmetic and checks zero remainders
before dividing $e_G$ by $|\Aut(G)|=2|G|=921631011840$.
Both old $M_{11}$ table names are retained and checked separately.
\begin{center}\begin{tabular}{rr@{\qquad}rr}\toprule
$m$&$n_G(m,11)$&$m$&$n_G(m,11)$\\\midrule
3&37&12&3583\\
4&503&14&4601\\
5&718&15&5814\\
6&1776&16&32488\\
7&2687&19&20493\\
8&8096&20&13094\\
10&6532&28&9202\\
11&11839&31&8325\\
\bottomrule\end{tabular}\end{center}

All sixteen numbers agree with the fixed published table~\cite{CL}.
The nine newly determined cells are
\[
 (3,11),(4,11),(5,11),(6,11),(7,11),(8,11),(10,11),(11,11),(11,15).
\]
Together with the previously accepted forty-six cells, including the
named Woldar exclusion, this gives fifty-five of 136 exact cells.
The remaining 81 cells, the full total and the chirality claims are
not settled by this component. The class-length subtraction is an
application of the existing generated-subgroup partition. The
classification data and character tables remain attributed premises;
no new subgroup classification or conjecture settlement is claimed.
\begin{thebibliography}{9}
\bibitem{CL} T. Connor and D. Leemans, \emph{Algorithmic enumeration of
regular maps}, Ars Math. Contemp. 10 (2016), 211--222.
\bibitem{atlas} T. Connor and D. Leemans, \emph{An atlas of subgroup
lattices of finite almost simple groups}, Ars Math. Contemp. 8 (2015),
259--266; retained online O'Nan table, rows 6, 23, 24, 119 and 286,
\url{https://leemans.dimitri.web.ulb.be/atlaslat/on.pdf}.
\end{thebibliography}
\end{document}
